\chapter{Testing and Review}
Log messages are part of your operational API. Tests keep them stable. Review keeps them readable.

\section{pytest caplog Basics}
\begin{lstlisting}[language=Python]
import logging

def test_log_message(caplog):
    logger = logging.getLogger("worker")
    with caplog.at_level(logging.INFO):
        logger.info(
            "queue.processed: outcome=success",
            extra={
                "event_name": "queue.processed",
                "outcome": "success",
                "ids": {"message_id": "m-9"},
                "fields": {"duration_ms": 12, "rows_out": 3},
            },
        )

    record = caplog.records[0]
    assert record.event_name == "queue.processed"
    assert record.outcome == "success"
\end{lstlisting}

\section{Checking Required Fields}
\begin{lstlisting}[language=Python]
import logging

def test_required_fields(caplog):
    logger = logging.getLogger("pipeline")
    with caplog.at_level(logging.WARNING):
        logger.warning(
            "batch.load: outcome=retry",
            extra={
                "event_name": "batch.load",
                "outcome": "retry",
                "fields": {"attempt": 2, "backoff_ms": 1000},
                "ids": {"job_id": "job_22"},
            },
        )

    record = caplog.records[0]
    assert "job_id" in record.ids
    assert record.fields["attempt"] == 2
\end{lstlisting}

\section{Checklist: Review and Tests}
\begin{itemize}
\item Tests assert event name and outcome.
\item Required IDs are present.
\item Fields include duration and attempt when relevant.
\item Levels match the outcome.
\end{itemize}

